\documentclass[11pt,a4paper]{article}

% ==================== TEMEL PAKETLER (xelatex) ====================
\usepackage{fontspec}
\usepackage{polyglossia}
\setmainlanguage{turkish}
\usepackage[a4paper,margin=2cm,top=2.4cm,bottom=2.4cm]{geometry}
\usepackage{xcolor}
\usepackage{graphicx}
\usepackage{array}
\usepackage{booktabs}
\usepackage{longtable}
\usepackage{tabularx}
\usepackage{colortbl}
\usepackage{multirow}
\usepackage{enumitem}
\usepackage[most]{tcolorbox}
\usepackage{listings}
\usepackage{fancyhdr}
\usepackage{titlesec}
\usepackage{titletoc}
\usepackage{hyperref}
\usepackage{fontawesome5}
\setlength{\headheight}{13.6pt}

% ==================== FONTLAR ====================
\setmainfont{Latin Modern Roman}
\newfontfamily\codefont{Consolas}[Scale=0.88]

% ==================== RENK PALETİ ====================
\definecolor{anaMavi}{HTML}{1B3A5B}
\definecolor{vurguMavi}{HTML}{2E6DA4}
\definecolor{acikMavi}{HTML}{EAF2F9}
\definecolor{turuncu}{HTML}{D9701E}
\definecolor{yesil}{HTML}{2E7D32}
\definecolor{acikYesil}{HTML}{E8F5E9}
\definecolor{kirmizi}{HTML}{C0392B}
\definecolor{mor}{HTML}{6C3483}
\definecolor{acikMor}{HTML}{F3E9F7}
\definecolor{acikGri}{HTML}{F4F5F7}
\definecolor{ortaGri}{HTML}{6B7280}
\definecolor{koyuGri}{HTML}{2D2D2D}
\definecolor{kodArka}{HTML}{FBFBFA}
\definecolor{satirRenk}{HTML}{F0F4F8}
\definecolor{qtYesil}{HTML}{41CD52}

% ==================== HYPERREF ====================
\hypersetup{
  colorlinks=true,
  linkcolor=vurguMavi,
  urlcolor=vurguMavi,
  pdftitle={QML Temelleri ve Pratik Örnekler Rehberi},
  pdfauthor={QML Rehberi},
  bookmarksnumbered=true
}

% ==================== BAŞLIK STİLLERİ ====================
\titleformat{\section}
  {\color{anaMavi}\Large\bfseries}
  {\thesection.}{0.6em}{}
  [\vspace{2pt}{\color{turuncu}\titlerule[1.6pt]}]
\titleformat{\subsection}
  {\color{vurguMavi}\large\bfseries}
  {\thesubsection}{0.6em}{}
\titleformat{\subsubsection}
  {\color{koyuGri}\normalsize\bfseries}
  {\thesubsubsection}{0.6em}{}
\titlespacing*{\section}{0pt}{16pt}{8pt}

% ==================== BAŞLIK / ALTBİLGİ ====================
\pagestyle{fancy}
\fancyhf{}
\renewcommand{\headrulewidth}{0.4pt}
\renewcommand{\footrulewidth}{0.4pt}
\fancyhead[L]{\small\color{ortaGri}QML Temelleri Rehberi}
\fancyhead[R]{\small\color{ortaGri}\leftmark}
\fancyfoot[C]{\small\color{ortaGri}\thepage}

% ==================== KOD LİSTELEME: QML ====================
\lstdefinelanguage{QML}{
  morekeywords={import,property,signal,function,id,anchors,width,height,color,
    Rectangle,Item,Text,Image,MouseArea,Component,Window,ApplicationWindow,Column,
    Row,Grid,Flow,RowLayout,ColumnLayout,GridLayout,Layout,ListView,GridView,
    ListModel,ListElement,Repeater,Button,TextField,TextArea,Slider,CheckBox,
    RadioButton,ComboBox,Connections,readonly,alias,required,default,var,int,real,
    string,bool,true,false,null,on,Behavior,states,State,PropertyChanges,Transition,
    Timer,Loader,pragma,QtQuick,Controls,Layouts,NumberAnimation,PropertyAnimation,
    SequentialAnimation,ParallelAnimation,ColorAnimation,RotationAnimation,
    onClicked,onCompleted,onPressed,onReleased,model,delegate,sourceComponent,
    TapHandler,DragHandler,Gradient,GradientStop,Flickable,ScrollView,Keys,
    when,anchors,parent,elide,wrapMode,font,source,fillMode,spacing,scale,rotation,
    opacity,visible,enabled,focus,z,radius,border},
  sensitive=true,
  morecomment=[l]{//},
  morecomment=[s]{/*}{*/},
  morestring=[b]",
  morestring=[b]'
}
\lstdefinestyle{qmlstyle}{
  language=QML,
  basicstyle=\codefont\small\color{koyuGri},
  keywordstyle=\color{mor}\bfseries,
  commentstyle=\color{yesil}\itshape,
  stringstyle=\color{turuncu},
  numberstyle=\tiny\color{ortaGri},
  numbers=left,
  numbersep=8pt,
  backgroundcolor=\color{qtYesil!6},
  frame=single,
  framesep=6pt,
  rulecolor=\color{qtYesil!50!ortaGri},
  breaklines=true,
  showstringspaces=false,
  tabsize=4,
  columns=fullflexible,
  keepspaces=true
}
\lstset{style=qmlstyle}

% ==================== ÖZEL KUTULAR ====================
\newtcolorbox{bilgikutu}[1][Bilgi]{
  colback=acikMavi, colframe=vurguMavi, boxrule=0.8pt, arc=3pt,
  left=8pt,right=8pt,top=6pt,bottom=6pt,
  fonttitle=\bfseries\color{white}, coltitle=white,
  title={\faInfoCircle~~#1}
}
\newtcolorbox{ipucukutu}[1][İpucu]{
  colback=acikYesil, colframe=yesil, boxrule=0.8pt, arc=3pt,
  left=8pt,right=8pt,top=6pt,bottom=6pt,
  fonttitle=\bfseries\color{white}, coltitle=white,
  title={\faLightbulb~~#1}
}
\newtcolorbox{uyarikutu}[1][Dikkat]{
  colback=kirmizi!8, colframe=kirmizi, boxrule=0.8pt, arc=3pt,
  left=8pt,right=8pt,top=6pt,bottom=6pt,
  fonttitle=\bfseries\color{white}, coltitle=white,
  title={\faExclamationTriangle~~#1}
}
\newtcolorbox{ornekkutu}[1][Örnek]{
  colback=acikGri, colframe=ortaGri!60, boxrule=0.7pt, arc=2pt,
  left=8pt,right=8pt,top=6pt,bottom=6pt,
  fonttitle=\bfseries\color{koyuGri}, coltitle=koyuGri,
  title={\faCode~~#1}
}
\newtcolorbox{notkutu}[1][Not]{
  colback=acikMor, colframe=mor, boxrule=0.8pt, arc=3pt,
  left=8pt,right=8pt,top=6pt,bottom=6pt,
  fonttitle=\bfseries\color{white}, coltitle=white,
  title={\faStickyNote~~#1}
}

% Yardımcı komutlar
\newcommand{\kls}[1]{\texttt{\color{anaMavi}\textbf{#1}}}
\newcommand{\hx}[1]{\texttt{\color{turuncu}#1}}
\newcommand{\kod}[1]{\texttt{\color{koyuGri}#1}}

\renewcommand{\arraystretch}{1.35}

% ==================== BAŞLANGIÇ ====================
\begin{document}

% ---------- KAPAK ----------
\begin{titlepage}
\centering
\vspace*{1.6cm}
{\color{qtYesil}\rule{\textwidth}{2pt}}\\[6pt]
{\color{anaMavi}\Huge\bfseries QML Temelleri}\\[10pt]
{\color{vurguMavi}\LARGE Pratik Örneklerle Qt Quick}\\[8pt]
{\color{ortaGri}\large Bileşenler \textbullet\ Layout'lar \textbullet\ Modeller \textbullet\ Animasyonlar}\\[6pt]
{\color{qtYesil}\rule{\textwidth}{2pt}}\\[1.6cm]

\begin{tcolorbox}[width=0.88\textwidth,colback=acikMavi,colframe=vurguMavi,arc=4pt,boxrule=1pt]
\centering\large
Bu belge, Qt Quick / QML ile arayüz geliştirmede \textbf{en sık kullanılan}
yapı taşlarını sıfırdan anlatır: görsel öğeler, özellik bağlama, konumlandırma,
\textbf{yeniden kullanılabilir bileşenlere bölme}, sinyaller, durumlar,
animasyonlar ve liste görünümleri. Her konu \textbf{çok sayıda kopyala-çalıştır
örneğiyle} verilmiştir.
\end{tcolorbox}

\vspace{0.9cm}
\begin{tcolorbox}[width=0.88\textwidth,colback=acikGri,colframe=ortaGri!50,arc=3pt,boxrule=0.6pt]
\small
\textbf{Kapsam:} Item/Rectangle/Text/Image \textbullet\ property \& binding
\textbullet\ anchors \textbullet\ Row/Column/Grid \& Layouts \textbullet\
\textbf{kendi bileşenini yazma} \textbullet\ property/signal API \textbullet\
sinyaller \& handler \textbullet\ States \& Transitions \textbullet\ animasyonlar
\& Behavior \textbullet\ ListModel/ListView/Repeater/GridView \textbullet\ Loader
\textbullet\ Qt Quick Controls \textbullet\ en iyi pratikler
\end{tcolorbox}

\vfill
{\color{ortaGri}\large Hazırlanma Tarihi: 9 Temmuz 2026}\\[4pt]
{\color{ortaGri} Qt 6.x \textbullet\ QtQuick \textbullet\ QtQuick.Controls \& Layouts}
\vspace*{0.6cm}
\end{titlepage}

% ---------- İÇİNDEKİLER ----------
\tableofcontents
\thispagestyle{empty}
\clearpage

% =====================================================================
\section{QML Söz Dizimi: Temel Kavramlar}
% =====================================================================

QML, kullanıcı arayüzünü \emph{ne olduğunu} tanımlayarak (deklaratif) yazdığınız
bir dildir. Bir QML dosyası, iç içe geçmiş \textbf{nesnelerden} oluşan bir ağaçtır.
Her nesnenin \emph{özellikleri} (property) vardır ve bu özellikler birbirine
\emph{bağlanabilir} (binding).

\subsection{İlk QML Dosyası}

\begin{lstlisting}
import QtQuick               // Temel QML tipleri (Item, Rectangle, Text...)
import QtQuick.Window        // Window tipi

Window {
    width: 400
    height: 300
    visible: true
    title: "Ilk QML Uygulamam"

    Rectangle {
        anchors.centerIn: parent   // Pencerenin ortasina yerlestir
        width: 200
        height: 100
        color: "steelblue"
        radius: 12                 // Yuvarlatilmis koseler

        Text {
            anchors.centerIn: parent
            text: "Merhaba QML"
            color: "white"
            font.pixelSize: 20
        }
    }
}
\end{lstlisting}

\subsection{Nesne, id ve Özellik}

Her QML nesnesi \kod{TipAdi \{ ... \}} biçiminde yazılır. İçine özellik atamaları
ve alt nesneler koyarsınız.

\begin{center}
\begin{tabularx}{\textwidth}{@{}>{\ttfamily\bfseries\color{vurguMavi}}l >{\raggedright\arraybackslash}X@{}}
\toprule
\rowcolor{acikMavi}\color{anaMavi}\textbf{Öğe} & \color{anaMavi}\textbf{Açıklama} \\
\midrule
id & Nesneye benzersiz kimlik; başka nesnelerden erişmek için. Örn: \kod{id: kutu} \\
\rowcolor{acikGri} property & Özellik: nesnenin verisi. \kod{width}, \kod{color}, \kod{text}... \\
parent & Bir nesnenin görsel ebeveyni; \kod{anchors} ve boyutlandırmada kullanılır. \\
\rowcolor{acikGri} \{ \} & Nesne kapsamı; içine alt nesneler ve özellikler yazılır. \\
\bottomrule
\end{tabularx}
\end{center}

\begin{ipucukutu}
\kod{id} bir string değildir, tırnak almaz: \kod{id: baslik}. Aynı dosyada
benzersiz olmalı, küçük harfle başlamalıdır. Bir nesneye başka yerden erişmenin
tek yolu \kod{id}'sidir: \kod{baslik.text = "Yeni"}.
\end{ipucukutu}

\subsection{Property Binding: QML'in Kalbi}

Bir özelliği sabit değil, \emph{bir ifadeye} bağlayabilirsiniz. İfadedeki değerler
değişince özellik \textbf{otomatik güncellenir}.

\begin{lstlisting}
Rectangle {
    id: kutu
    width: 200
    // height artik width'e bagli; width degisince height da degisir
    height: width / 2
    color: width > 150 ? "green" : "red"   // Kosula bagli renk
}
\end{lstlisting}

\begin{bilgikutu}
Binding tek yönlü ve canlıdır. \kod{height: width / 2} yazdığınızda, "şu an
width'in yarısını al" değil, "\emph{her zaman} width'in yarısı ol" dersiniz.
\kod{width} 300 olursa \kod{height} anında 150 olur. Elle güncelleme kodu
yazmazsınız — bu, QML'i güçlü kılan özelliktir.
\end{bilgikutu}

% =====================================================================
\section{Temel Görsel Öğeler}
% =====================================================================

QtQuick'te en sık kullanacağınız birkaç temel tip vardır. Bunları öğrenmek
arayüzlerin \%80'ini kapsar.

\subsection{Item — Görünmez Kapsayıcı}

\kls{Item} tüm görsel öğelerin taban tipidir; kendi başına bir şey çizmez ama
gruplamak, konumlandırmak ve boyutlandırmak için kullanılır.

\begin{lstlisting}
Item {
    width: 300; height: 200
    // Icindeki ogeleri gruplar; kendisi gorunmez
    Rectangle { anchors.fill: parent; color: "lightgray" }
}
\end{lstlisting}

\subsection{Rectangle — Dikdörtgen}

\begin{lstlisting}
Rectangle {
    width: 120; height: 80
    color: "tomato"
    radius: 10                    // Kose yuvarlakligi
    border.color: "darkred"       // Kenarlik rengi
    border.width: 2               // Kenarlik kalinligi
    opacity: 0.9                  // Saydamlik (0=gorunmez, 1=opak)

    // Renk gecisi (gradient)
    gradient: Gradient {
        GradientStop { position: 0.0; color: "orange" }
        GradientStop { position: 1.0; color: "red" }
    }
}
\end{lstlisting}

\subsection{Text — Metin}

\begin{lstlisting}
Text {
    text: "Uzun bir metin bu satira sigmayabilir"
    color: "#333333"
    font.pixelSize: 18
    font.bold: true
    font.family: "Arial"
    width: 150                    // Sabit genislik
    wrapMode: Text.WordWrap       // Kelime bazli alt satira gec
    elide: Text.ElideRight        // Sigmazsa sonuna "..." koy
    horizontalAlignment: Text.AlignHCenter
}
\end{lstlisting}

\subsection{Image — Görsel}

\begin{lstlisting}
Image {
    source: "qrc:/images/logo.png"   // veya "https://..." / dosya yolu
    width: 100; height: 100
    fillMode: Image.PreserveAspectFit  // Orani koru, sigdir
    asynchronous: true                 // Arka planda yukle (UI donmasin)
}
\end{lstlisting}

\begin{center}
\begin{tabularx}{\textwidth}{@{}>{\ttfamily\color{anaMavi}\small}l >{\raggedright\arraybackslash}X@{}}
\toprule
\rowcolor{acikMavi}\color{anaMavi}\textbf{fillMode} & \color{anaMavi}\textbf{Davranış} \\
\midrule
Image.Stretch & Görseli çerçeveye uzat (oran bozulur) \\
\rowcolor{acikGri} Image.PreserveAspectFit & Oranı koruyarak sığdır (boşluk kalabilir) \\
Image.PreserveAspectCrop & Oranı koruyarak doldur (taşan kırpılır) \\
\rowcolor{acikGri} Image.Tile & Döşe (tekrarla) \\
\bottomrule
\end{tabularx}
\end{center}

\subsection{MouseArea — Fare/Dokunma Etkileşimi}

Görsel öğeler kendiliğinden tıklanabilir değildir; üzerlerine görünmez bir
\kls{MouseArea} koyarsınız.

\begin{lstlisting}
Rectangle {
    id: dugme
    width: 120; height: 40
    color: fare.pressed ? "darkblue" : "steelblue"   // Basiliyken koyulasir
    radius: 6

    Text { anchors.centerIn: parent; text: "Tikla"; color: "white" }

    MouseArea {
        id: fare
        anchors.fill: parent          // Tum dikdortgeni kapla
        hoverEnabled: true            // Fare uzerine gelince de algila
        onClicked: console.log("Tiklandi!")
        onDoubleClicked: console.log("Cift tiklandi")
        onPressed: console.log("Basildi")
        onEntered: dugme.scale = 1.1  // Fare geldi
        onExited:  dugme.scale = 1.0  // Fare cikti
    }
}
\end{lstlisting}

\begin{notkutu}[Modern alternatif: Handler'lar]
Qt 6'da \kls{MouseArea} yerine daha esnek \textbf{Input Handler}'lar önerilir:
\kod{TapHandler} (tıklama), \kod{DragHandler} (sürükleme), \kod{HoverHandler}
(fare üzeri), \kod{PinchHandler} (kıstırma). Bunlar birden fazla aynı anda
çalışabilir ve dokunmatik ekranlarla daha uyumludur.
\end{notkutu}

\begin{lstlisting}
Rectangle {
    width: 100; height: 100; color: "coral"
    TapHandler { onTapped: console.log("TapHandler ile tiklandi") }
    DragHandler { }   // Tek satirla surukle-birak!
}
\end{lstlisting}

% =====================================================================
\section{Konumlandırma: Anchors ve Positioner'lar}
% =====================================================================

QML'de öğeleri yerleştirmenin üç yolu vardır: sabit x/y, \textbf{anchors}
(bağlama), \textbf{positioner}'lar (Row/Column/Grid) ve \textbf{Layout}'lar.

\subsection{Anchors — Bağlama}

Bir öğenin kenarlarını başka bir öğenin kenarlarına "yapıştırırsınız".

\begin{lstlisting}
Rectangle {
    id: kapsayici
    width: 300; height: 200; color: "#eee"

    Rectangle {
        width: 80; height: 80; color: "teal"
        anchors.top: parent.top          // Ust kenari ebeveynin ustune
        anchors.left: parent.left        // Sol kenari ebeveynin soluna
        anchors.margins: 10              // 10px bosluk birak
    }
    Rectangle {
        width: 80; height: 80; color: "orange"
        anchors.centerIn: parent         // Tam ortala
    }
    Rectangle {
        height: 40; color: "purple"
        anchors.left: parent.left
        anchors.right: parent.right      // Sol+sag -> genislik otomatik
        anchors.bottom: parent.bottom
    }
}
\end{lstlisting}

\begin{center}
\begin{tabularx}{\textwidth}{@{}>{\ttfamily\color{anaMavi}\small}l >{\raggedright\arraybackslash}X@{}}
\toprule
\rowcolor{acikMavi}\color{anaMavi}\textbf{Anchor} & \color{anaMavi}\textbf{Açıklama} \\
\midrule
anchors.fill: parent & Ebeveyni tamamen kapla \\
\rowcolor{acikGri} anchors.centerIn: parent & Yatay+dikey ortala \\
anchors.top / bottom / left / right & Tek kenar bağlama \\
\rowcolor{acikGri} anchors.horizontalCenter & Yatay merkez hizalama \\
anchors.margins / topMargin ... & Bağlanan kenardan boşluk \\
\bottomrule
\end{tabularx}
\end{center}

\begin{uyarikutu}
Bir öğeyi hem \kod{anchors} ile hem de \kod{x}/\kod{y} ile konumlandırmayın —
çakışırlar. Ayrıca \kod{anchors} yalnızca \textbf{ebeveyn} veya \textbf{kardeş}
öğelere bağlanabilir, rastgele bir öğeye değil.
\end{uyarikutu}

\subsection{Row, Column, Grid — Positioner'lar}

Öğeleri otomatik olarak yatay, dikey veya ızgara biçiminde dizerler.

\begin{lstlisting}
Column {                       // Alt alta dizer
    spacing: 10                // Ogeler arasi bosluk
    Rectangle { width: 100; height: 30; color: "red" }
    Rectangle { width: 100; height: 30; color: "green" }
    Rectangle { width: 100; height: 30; color: "blue" }
}

Row {                          // Yan yana dizer
    spacing: 8
    Repeater {                 // 5 tane uret
        model: 5
        Rectangle { width: 40; height: 40; color: "steelblue"; radius: 4 }
    }
}

Grid {                         // Izgara: 3 sutun
    columns: 3
    spacing: 6
    Repeater {
        model: 9
        Rectangle { width: 50; height: 50; color: "orange" }
    }
}
\end{lstlisting}

\subsection{Layout'lar — Esnek ve Uyarlanabilir}

\kls{RowLayout}, \kls{ColumnLayout}, \kls{GridLayout}; öğelerin pencere boyutuna
göre \emph{esnemesini} sağlar. Bunun için \kod{QtQuick.Layouts} import edilir ve
öğeler \kod{Layout.*} ekli özellikleriyle yapılandırılır.

\begin{lstlisting}
import QtQuick
import QtQuick.Layouts

RowLayout {
    anchors.fill: parent
    spacing: 10

    Rectangle {
        color: "salmon"
        Layout.preferredWidth: 100
        Layout.fillHeight: true       // Dikeyde tum alani doldur
    }
    Rectangle {
        color: "khaki"
        Layout.fillWidth: true        // Kalan yatay alani KAP
        Layout.fillHeight: true
    }
    Rectangle {
        color: "lightgreen"
        Layout.preferredWidth: 60
        Layout.fillHeight: true
    }
}
\end{lstlisting}

\begin{center}
\begin{tabularx}{\textwidth}{@{}>{\ttfamily\color{anaMavi}\small}l >{\raggedright\arraybackslash}X@{}}
\toprule
\rowcolor{acikMavi}\color{anaMavi}\textbf{Layout özelliği} & \color{anaMavi}\textbf{Açıklama} \\
\midrule
Layout.fillWidth / fillHeight & Kalan alanı doldur (esne) \\
\rowcolor{acikGri} Layout.preferredWidth / Height & Tercih edilen boyut \\
Layout.minimumWidth / maximumWidth & Alt/üst boyut sınırı \\
\rowcolor{acikGri} Layout.alignment & Hücre içi hizalama (Qt.AlignCenter...) \\
Layout.columnSpan / rowSpan & GridLayout'ta hücre yayılımı \\
\bottomrule
\end{tabularx}
\end{center}

\begin{ipucukutu}[Positioner mı, Layout mu?]
\kls{Row/Column} basit, sabit boyutlu dizilimler için hafif ve hızlıdır.
\kls{RowLayout/ColumnLayout} ise öğelerin \textbf{pencere boyutuna göre esnemesi}
gerektiğinde (responsive arayüz) kullanılır. \kod{Layout.fillWidth} gibi
özellikler sadece Layout'lar içinde çalışır, Row/Column içinde çalışmaz.
\end{ipucukutu}

% =====================================================================
\section{Bileşenlere Bölme: Yeniden Kullanılabilir QML}
% =====================================================================

QML'in en güçlü tarafı: \textbf{her \kod{.qml} dosyası aynı zamanda bir tiptir}.
\kod{MyButton.qml} adında bir dosya oluşturursanız, başka QML dosyalarında
\kod{MyButton \{ \}} yazarak onu bir öğe gibi kullanabilirsiniz. Böylece tekrar
eden arayüz parçalarını bileşenlere bölersiniz.

\subsection{İlk Özel Bileşen}

Diyelim aynı stildeki düğmeyi defalarca kullanacaksınız. Ayrı bir dosyaya
taşıyın:

\begin{ornekkutu}[StyledButton.qml — kendi bileşenimiz]
\end{ornekkutu}
\begin{lstlisting}
import QtQuick

Rectangle {                        // Kok oge = bilesenin "govdesi"
    id: root
    width: 140; height: 44
    radius: 8
    color: area.pressed ? Qt.darker(baseColor, 1.2)
                        : (area.containsMouse ? Qt.lighter(baseColor, 1.1)
                                              : baseColor)

    // ---- DISARIYA ACILAN API ----
    property string text: "Dugme"       // Ozellestirilebilir etiket
    property color  baseColor: "#2E6DA4"
    signal clicked()                    // Disariya bildirilen sinyal

    Text {
        anchors.centerIn: parent
        text: root.text                 // Disaridan gelen metin
        color: "white"
        font.pixelSize: 16
    }

    MouseArea {
        id: area
        anchors.fill: parent
        hoverEnabled: true
        onClicked: root.clicked()       // Kendi sinyalimizi yay
    }
}
\end{lstlisting}

Artık başka dosyada kullanabiliriz (aynı klasördeyse \kod{import} bile gerekmez):

\begin{lstlisting}
import QtQuick
import QtQuick.Window

Window {
    width: 300; height: 240; visible: true

    Column {
        anchors.centerIn: parent
        spacing: 12

        StyledButton {
            text: "Kaydet"
            baseColor: "#2E7D32"
            onClicked: console.log("Kaydedildi")   // Sinyali dinle
        }
        StyledButton {
            text: "Sil"
            baseColor: "#C0392B"
            onClicked: console.log("Silindi")
        }
        StyledButton { text: "Iptal" }             // Varsayilan mavi
    }
}
\end{lstlisting}

\begin{bilgikutu}
Bir bileşenin \textbf{genel API'si} üç şeyden oluşur: \textbf{(1)} dışarıdan
ayarlanabilen \kod{property}'ler (girdi), \textbf{(2)} dışarıya olay bildiren
\kod{signal}'ler (çıktı), \textbf{(3)} çağrılabilen \kod{function}'lar. İç yapı
(MouseArea, Text vs.) gizli kalır — kullanıcı sadece bu API'yi görür. Bu,
kapsülleme (encapsulation) demektir.
\end{bilgikutu}

\subsection{property alias — İç Öğeyi Dışarı Açma}

Bazen iç öğenin bir özelliğine dışarıdan doğrudan erişim vermek istersiniz.
\kod{alias} bunu sağlar:

\begin{lstlisting}
// LabeledField.qml
import QtQuick
import QtQuick.Controls

Row {
    spacing: 8
    property alias label: etiket.text     // Ic Text'in text'ini disa ac
    property alias value: girdi.text       // Ic TextField'in text'ini ac

    Text { id: etiket; text: "Ad:" }
    TextField { id: girdi; placeholderText: "..." }
}
\end{lstlisting}

\begin{lstlisting}
LabeledField {
    label: "E-posta:"
    value: "ornek@mail.com"
    onValueChanged: console.log("Deger:", value)   // alias otomatik sinyal uretir
}
\end{lstlisting}

\begin{notkutu}[property vs alias]
\kod{property string x: "..."} \emph{yeni} bir değer saklar. \kod{property alias x:
oge.y} ise mevcut bir öğenin özelliğine \emph{takma ad} verir — okuma ve yazma
doğrudan o öğeye gider. Alias, iç bir bileşenin özelliğini/nesnesini API'ye dahil
etmenin temiz yoludur.
\end{notkutu}

\subsection{default property — İçerik Yuvası}

Kendi bileşeninize, tıpkı \kls{Column} gibi, doğrudan çocuk öğe alabilme
yeteneği verebilirsiniz:

\begin{lstlisting}
// Card.qml -- icine konan her sey "content" alanina gider
import QtQuick

Rectangle {
    id: root
    radius: 10; color: "white"; border.color: "#ddd"
    implicitWidth: content.implicitWidth + 24
    implicitHeight: content.implicitHeight + 24

    default property alias data: content.data   // <-- Sihir burada

    Column {
        id: content
        anchors.centerIn: parent
        spacing: 8
    }
}
\end{lstlisting}

\begin{lstlisting}
Card {
    // Bunlar dogrudan Card'in icine (content Column'a) gider:
    Text { text: "Baslik"; font.bold: true }
    Text { text: "Aciklama metni burada" }
}
\end{lstlisting}

\subsection{inline Component ve Component Tipi}

Küçük, sadece o dosyada kullanılacak bileşenler için ayrı dosya açmadan
\kod{component} anahtar kelimesiyle satır içi tip tanımlayabilirsiniz:

\begin{lstlisting}
Item {
    // Dosya ici bilesen tanimi (Qt 6)
    component Rozet: Rectangle {
        width: 60; height: 24; radius: 12; color: "purple"
        property alias text: t.text
        Text { id: t; anchors.centerIn: parent; color: "white" }
    }

    Row {
        spacing: 6
        Rozet { text: "Yeni" }
        Rozet { text: "Pro"; color: "orange" }
    }
}
\end{lstlisting}

% =====================================================================
\section{Sinyaller ve Handler'lar}
% =====================================================================

Sinyaller, "bir şey oldu" bilgisini iletir. Yerleşik öğelerin hazır sinyalleri
(\kod{onClicked}, \kod{onPressed}...) vardır; kendi sinyallerinizi de
tanımlayabilirsiniz.

\subsection{Kendi Sinyalini Tanımlama ve Yayma}

\begin{lstlisting}
Item {
    id: kutu

    // Parametreli sinyal tanimi (Qt 6 tarzi)
    signal itemSelected(int index, string name)

    function sec(i, ad) {
        kutu.itemSelected(i, ad)      // Sinyali yay
    }

    // Ayni Item icinde handler:
    onItemSelected: (index, name) => {
        console.log("Secilen:", index, name)
    }
}
\end{lstlisting}

\subsection{Connections ile Dışarıdan Bağlanma}

Bir öğenin sinyaline, tanımından \emph{başka} bir yerde bağlanmak için:

\begin{lstlisting}
Timer {
    id: sayac
    interval: 1000; running: true; repeat: true
}

Connections {
    target: sayac
    function onTriggered() {           // Timer'in "triggered" sinyali
        console.log("Bir saniye gecti")
    }
}
\end{lstlisting}

\subsection{Property Değişim Sinyalleri}

Her özellik otomatik olarak bir \kod{on\textless Ozellik\textgreater Changed}
handler'ı sağlar:

\begin{lstlisting}
Rectangle {
    width: 100; height: 100
    // width degistiginde otomatik cagrilir
    onWidthChanged: console.log("Yeni genislik:", width)
    color: "gray"
    onColorChanged: console.log("Renk degisti")
}
\end{lstlisting}

% =====================================================================
\section{Durumlar (States) ve Geçişler (Transitions)}
% =====================================================================

Bir öğenin farklı görünümlerini (açık/kapalı, seçili/seçilmemiş) \textbf{state}
olarak tanımlar, aralarındaki \emph{yumuşak geçişi} \textbf{transition} ile
sağlarsınız.

\subsection{States ile Görünüm Değiştirme}

\begin{lstlisting}
Rectangle {
    id: panel
    width: 200; height: 60
    color: "steelblue"

    // Baslangicta state yok (varsayilan); tanimli durumlar:
    states: [
        State {
            name: "genisletilmis"
            PropertyChanges { target: panel; height: 200; color: "seagreen" }
        },
        State {
            name: "gizli"
            PropertyChanges { target: panel; height: 0; opacity: 0 }
        }
    ]

    MouseArea {
        anchors.fill: parent
        // Tikla -> durumu degistir (bosluk = varsayilana don)
        onClicked: panel.state = (panel.state === "genisletilmis")
                                 ? "" : "genisletilmis"
    }
}
\end{lstlisting}

\subsection{Transitions ile Yumuşak Geçiş}

State değişimi ani olur; \kls{Transition} eklerseniz animasyonlu olur:

\begin{lstlisting}
Rectangle {
    id: kutu
    width: 100; height: 100; color: "coral"
    states: State {
        name: "buyuk"
        PropertyChanges { target: kutu; width: 200; height: 200 }
    }
    transitions: Transition {
        // Tum sayisal ozellikleri 300ms'de yumusakca degistir
        NumberAnimation {
            properties: "width,height"
            duration: 300
            easing.type: Easing.InOutQuad
        }
    }
    MouseArea { anchors.fill: parent
                onClicked: kutu.state = kutu.state === "buyuk" ? "" : "buyuk" }
}
\end{lstlisting}

% =====================================================================
\section{Animasyonlar}
% =====================================================================

QML animasyonlar için zengin bir sistem sunar. En sık kullanılanlar:

\subsection{Behavior — Otomatik Yumuşatma}

Bir özellik \emph{her ne zaman} değişirse değişsin animasyonlu olsun istiyorsanız
\kls{Behavior} en pratik yoldur:

\begin{lstlisting}
Rectangle {
    id: top
    width: 60; height: 60; radius: 30; color: "tomato"
    x: 0

    // x her degistiginde 400ms'de yumusakca kaysin
    Behavior on x {
        NumberAnimation { duration: 400; easing.type: Easing.OutBounce }
    }

    MouseArea {
        anchors.fill: parent
        onClicked: top.x = top.x === 0 ? 200 : 0   // x'i degistir, animasyon otomatik
    }
}
\end{lstlisting}

\subsection{Sıralı ve Paralel Animasyonlar}

\begin{lstlisting}
Rectangle {
    id: kare
    width: 80; height: 80; color: "mediumpurple"

    // Tiklaninca calisacak animasyon zinciri
    SequentialAnimation {
        id: anim
        NumberAnimation { target: kare; property: "rotation"
                          to: 360; duration: 500 }
        ParallelAnimation {                        // Ayni anda:
            NumberAnimation { target: kare; property: "scale"
                              to: 1.5; duration: 300 }
            ColorAnimation  { target: kare; property: "color"
                              to: "gold"; duration: 300 }
        }
        NumberAnimation { target: kare; property: "scale"
                          to: 1.0; duration: 200 }
    }
    MouseArea { anchors.fill: parent; onClicked: anim.start() }
}
\end{lstlisting}

\begin{center}
\begin{tabularx}{\textwidth}{@{}>{\ttfamily\color{anaMavi}\small}l >{\raggedright\arraybackslash}X@{}}
\toprule
\rowcolor{acikMavi}\color{anaMavi}\textbf{Animasyon tipi} & \color{anaMavi}\textbf{Kullanım} \\
\midrule
NumberAnimation & Sayısal özellikler (x, width, opacity...) \\
\rowcolor{acikGri} ColorAnimation & Renk geçişleri \\
RotationAnimation & Dönüş (en kısa yol seçebilir) \\
\rowcolor{acikGri} SequentialAnimation & İçindekileri sırayla çalıştırır \\
ParallelAnimation & İçindekileri aynı anda çalıştırır \\
\rowcolor{acikGri} Behavior on X & X değişince otomatik animasyon \\
PropertyAnimation & Genel; her tip özelliğe uygulanır \\
\bottomrule
\end{tabularx}
\end{center}

\begin{ipucukutu}[Easing — yumuşama eğrileri]
\kod{easing.type} animasyonun "hissini" belirler: \kod{Easing.Linear} (sabit hız),
\kod{Easing.InOutQuad} (yavaş başla-hızlan-yavaşla), \kod{Easing.OutBounce}
(zıplayarak dur), \kod{Easing.OutElastic} (yaylanarak). Doğal his için genelde
\kod{InOutQuad} veya \kod{OutCubic} iyi başlangıçtır.
\end{ipucukutu}

% =====================================================================
\section{Modeller ve Görünümler: Liste Göstermek}
% =====================================================================

Tekrar eden veriyi (liste, ızgara) göstermenin yolu \textbf{model-view}
ayrımıdır: \emph{model} veriyi tutar, \emph{delegate} her öğenin nasıl
görüneceğini tanımlar, \emph{view} (ListView/GridView) ikisini birleştirir.

\subsection{Repeater — En Basit Tekrarlama}

\begin{lstlisting}
Column {
    spacing: 4
    Repeater {
        model: ["Elma", "Armut", "Muz"]   // Dizi model
        delegate: Text {
            required property string modelData   // Dizi elemani
            required property int index
            text: (index + 1) + ". " + modelData  // "1. Elma"
        }
    }
}
\end{lstlisting}

\subsection{ListModel + ListView — Kaydırılabilir Liste}

\begin{lstlisting}
import QtQuick

ListView {
    width: 250; height: 300
    spacing: 6
    clip: true                       // Tasan icerigi kirp

    model: ListModel {
        ListElement { ad: "Ali";  yas: 25 }
        ListElement { ad: "Ayse"; yas: 30 }
        ListElement { ad: "Veli"; yas: 28 }
    }

    delegate: Rectangle {
        // Model rolleri delegate icinde dogrudan erisilir
        required property string ad
        required property int yas
        width: ListView.view.width
        height: 50
        color: "#f0f4f8"
        radius: 6

        Row {
            anchors.verticalCenter: parent.verticalCenter
            anchors.left: parent.left; anchors.leftMargin: 12
            spacing: 12
            Text { text: ad; font.bold: true }
            Text { text: yas + " yas"; color: "gray" }
        }
    }
}
\end{lstlisting}

\begin{uyarikutu}
Delegate içinde model rollerine erişmek için Qt 6'da \kod{required property}
kullanımı \textbf{önerilir} (\kod{required property string ad}). Eski kodda bu
roller sihirli biçimde kapsamdaydı ama isim çakışması ve karışıklık yaratıyordu.
\kod{required} ile hangi rollerin kullanıldığı açık ve güvenli olur.
\end{uyarikutu}

\subsection{ListModel'i Dinamik Değiştirme}

\begin{lstlisting}
ListModel { id: liste }

// Ekle / sil / guncelle
function ekle(ad) { liste.append({ ad: ad, yas: 0 }) }
function sil(i)   { liste.remove(i) }
function temizle(){ liste.clear() }
function guncelle(i, yeniAd) { liste.setProperty(i, "ad", yeniAd) }
\end{lstlisting}

\subsection{GridView — Izgara Görünüm}

\begin{lstlisting}
GridView {
    width: 320; height: 300
    cellWidth: 100; cellHeight: 100     // Her hucrenin boyutu
    model: 12
    delegate: Rectangle {
        required property int index
        width: 90; height: 90
        color: Qt.hsla(index / 12, 0.6, 0.6, 1)   // Renk carki
        radius: 8
        Text { anchors.centerIn: parent; text: index; color: "white" }
    }
}
\end{lstlisting}

\begin{center}
\begin{tabularx}{\textwidth}{@{}>{\ttfamily\color{anaMavi}\small}l >{\raggedright\arraybackslash}X@{}}
\toprule
\rowcolor{acikMavi}\color{anaMavi}\textbf{Özellik} & \color{anaMavi}\textbf{Açıklama} \\
\midrule
model & Veri kaynağı (sayı, dizi, ListModel, C++ modeli) \\
\rowcolor{acikGri} delegate & Her öğenin görsel şablonu \\
clip: true & Görünüm dışına taşan öğeleri gizle \\
\rowcolor{acikGri} spacing & Öğeler arası boşluk (ListView) \\
header / footer & Liste başı/sonu öğeleri \\
\rowcolor{acikGri} ScrollBar.vertical & Kaydırma çubuğu ekleme \\
\bottomrule
\end{tabularx}
\end{center}

% =====================================================================
\section{Loader ve Dinamik İçerik}
% =====================================================================

\kls{Loader}, bir bileşeni \emph{gerektiğinde} yükler — sekmeli arayüzlerde,
koşullu içerikte ve belleği verimli kullanmada idealdir.

\begin{lstlisting}
Item {
    Loader {
        id: yukleyici
        anchors.fill: parent
        // Kaynagi degistirerek icerigi degistir
        source: durum === "ayarlar" ? "SettingsPage.qml" : "HomePage.qml"
    }

    // veya satir ici bilesen yukle:
    Loader {
        sourceComponent: acikMi ? detayBileseni : null
    }
    Component {
        id: detayBileseni
        Rectangle { color: "lightyellow"; width: 100; height: 100 }
    }
}
\end{lstlisting}

\begin{ipucukutu}
\kod{Loader.active: false} yaparak yüklenen içeriği bellekten kaldırabilirsiniz.
Ağır sayfaları yalnızca görünür olduklarında yüklemek (lazy loading) uygulama
başlangıcını hızlandırır. \kod{asynchronous: true} ile yükleme arka planda olur.
\end{ipucukutu}

% =====================================================================
\section{Qt Quick Controls: Hazır Bileşenler}
% =====================================================================

\kod{QtQuick.Controls}; düğme, metin kutusu, kaydırıcı gibi hazır, stillenebilir
arayüz bileşenleri sunar. Sıfırdan yazmak yerine bunları kullanın.

\begin{lstlisting}
import QtQuick
import QtQuick.Controls
import QtQuick.Layouts

ColumnLayout {
    anchors.centerIn: parent
    spacing: 10

    Button {
        text: "Tikla"
        onClicked: durumEtiketi.text = "Butona basildi"
    }

    TextField {
        placeholderText: "Adinizi girin"
        onTextChanged: console.log("Metin:", text)
    }

    Slider {
        from: 0; to: 100; value: 50
        onValueChanged: console.log("Deger:", value.toFixed(0))
    }

    CheckBox {
        text: "Sartlari kabul ediyorum"
        onCheckedChanged: console.log("Onay:", checked)
    }

    ComboBox {
        model: ["Kucuk", "Orta", "Buyuk"]
        onCurrentTextChanged: console.log("Secim:", currentText)
    }

    Label { id: durumEtiketi; text: "Hazir" }
}
\end{lstlisting}

\begin{center}
\begin{tabularx}{\textwidth}{@{}>{\ttfamily\color{anaMavi}\small}l >{\raggedright\arraybackslash}X@{}}
\toprule
\rowcolor{acikMavi}\color{anaMavi}\textbf{Control} & \color{anaMavi}\textbf{Kullanım} \\
\midrule
Button & Tıklanabilir düğme (\kod{onClicked}) \\
\rowcolor{acikGri} TextField / TextArea & Tek/çok satır metin girişi \\
Slider / RangeSlider & Kaydırıcı ile değer seçimi \\
\rowcolor{acikGri} CheckBox / Switch & Aç/kapa seçimi \\
RadioButton & Gruptan tek seçim \\
\rowcolor{acikGri} ComboBox & Açılır liste \\
Dialog / Popup & İletişim kutusu / açılır pencere \\
\rowcolor{acikGri} TabBar / StackView & Sekme ve sayfa yığını navigasyonu \\
\bottomrule
\end{tabularx}
\end{center}

\begin{notkutu}[Stil seçimi]
Qt Quick Controls farklı stillerle gelir: \kod{Basic}, \kod{Material},
\kod{Universal}, \kod{Fusion}, \kod{macOS}, \kod{Windows}. Uygulama genelinde
stili \kod{main.cpp}'de \kod{QQuickStyle::setStyle("Material")} ile veya
\kod{qtquickcontrols2.conf} dosyasıyla ayarlarsınız. Material stil, renk ve
tema özelleştirmesi için en esnek olanıdır.
\end{notkutu}

% =====================================================================
\section{En İyi Pratikler ve Sık Hatalar}
% =====================================================================

\subsection{Yapılması Gerekenler}

\begin{itemize}[leftmargin=1.4em]
\item \textbf{Bileşenlere böl:} 100+ satırlık QML dosyalarını mantıklı,
      yeniden kullanılabilir bileşenlere ayırın. Her bileşen tek bir işi yapsın.
\item \textbf{Net API tasarla:} Bileşenin dışarıya açtığı \kod{property} ve
      \kod{signal}'leri bilinçli seçin; iç yapıyı gizleyin.
\item \textbf{Binding kullan, elle güncelleme yapma:} \kod{width: parent.width /
      2} yaz; \kod{onParentWidthChanged} ile elle atama yapma.
\item \textbf{Layout'larla responsive tasarla:} Sabit piksel yerine
      \kod{anchors} ve \kod{Layout} kullan; farklı ekranlara uyum sağlasın.
\item \textbf{required property kullan:} Delegate'lerde model rollerini açıkça
      belirt.
\item \textbf{Ağır işi QML'de yapma:} Karmaşık hesaplama/döngüleri C++'a taşı;
      QML sadece sunum olsun.
\end{itemize}

\subsection{Kaçınılması Gerekenler}

\begin{center}
\begin{tabularx}{\textwidth}{@{}>{\raggedright\arraybackslash\bfseries\color{kirmizi}}p{5.2cm} >{\raggedright\arraybackslash}X@{}}
\toprule
\rowcolor{kirmizi!12}\color{kirmizi}\textbf{Hata} & \color{kirmizi}\textbf{Neden kaçınmalı} \\
\midrule
Binding'i JS ile ezmek & \kod{width = 100} atarsanız binding \textbf{kalıcı olarak kırılır}. İfade tabanlı kalın. \\
\rowcolor{acikGri} anchors + x/y karıştırmak & Çakışır, öngörülemez sonuç. Birini seçin. \\
Dev tek dosya & Bakımı imkânsızlaşır. Bileşenlere bölün. \\
\rowcolor{acikGri} Delegate'te ağır iş & Kaydırma takılır. Delegate hafif olmalı; \kod{clip} ve \kod{asynchronous} kullanın. \\
id yerine sabit yol & \kod{parent.parent.parent...} kırılgan. \kod{id} ile referans verin. \\
\rowcolor{acikGri} console.log unutmak & Çalışmayan binding'i sessizce kaçırırsınız. Şüphede logla. \\
\bottomrule
\end{tabularx}
\end{center}

\subsection{Hızlı Referans}

\begin{center}
\begin{tabularx}{\textwidth}{@{}>{\ttfamily\color{anaMavi}\small}l >{\raggedright\arraybackslash}X@{}}
\toprule
\rowcolor{acikMavi}\color{anaMavi}\textbf{Öğe} & \color{anaMavi}\textbf{Amaç} \\
\midrule
Item / Rectangle / Text / Image & Temel görsel yapı taşları \\
\rowcolor{acikGri} MouseArea / TapHandler & Kullanıcı etkileşimi \\
anchors & Öğeyi ebeveyn/kardeşe bağlama \\
\rowcolor{acikGri} Row / Column / Grid & Otomatik dizilim \\
RowLayout / ColumnLayout & Esneyen (responsive) yerleşim \\
\rowcolor{acikGri} property / alias / signal & Bileşen API'si \\
States / Transitions & Görünüm durumları ve geçişleri \\
\rowcolor{acikGri} Behavior / NumberAnimation & Animasyon \\
ListView / GridView / Repeater & Liste ve ızgara görünümleri \\
\rowcolor{acikGri} Loader & Dinamik/tembel yükleme \\
Qt Quick Controls & Hazır arayüz bileşenleri \\
\bottomrule
\end{tabularx}
\end{center}

\vspace{0.6cm}
\begin{tcolorbox}[colback=qtYesil!8,colframe=qtYesil!70!koyuGri,arc=3pt,boxrule=1pt]
\centering
\textbf{Özet:} QML deklaratiftir — \emph{ne} istediğinizi yazarsınız, \emph{nasıl}
güncelleneceğini binding halleder. Arayüzü küçük, net API'li bileşenlere bölün;
yerleşimi \kod{anchors} ve Layout'larla responsive yapın; tekrarlı veriyi
model-view ile gösterin; hareketi States/Behavior ile canlandırın. İnce bir QML
katmanı + ayrı bileşenler = bakımı kolay, ölçeklenebilir arayüz.
\end{tcolorbox}

\end{document}